\section{Introduction}
An Agent receives the instruction \emph{process the next support ticket}. It reads the ticket, checks an order and approval, consults a model, and submits a refund. Encryption hides message contents from a network observer while contacted services, sizes and timing remain visible. The tool's answer then enters a model request, creating another observable representation. Adding cover raises a second question: does it end at a proxy, execute a read, or change business state? Figure~\ref{fig:chain} follows information and actions through this example.

\begin{figure*}[t]
\centering
\begin{tikzpicture}[x=1cm,y=1cm,font=\small,>=Latex]
\tikzset{box/.style={draw,rounded corners,align=center,minimum height=.85cm,text width=2.4cm},
 note/.style={align=center,font=\footnotesize,text width=3.1cm}}
\node[box,fill=blue!5] (u) at (1.3,1.15) {User\\``Process the next ticket''};
\node[box,fill=blue!5] (a) at (5.1,1.15) {Agent\\plan, check, act, retry};
\node[box,fill=green!6] (m) at (9.2,2.6) {Model service\\next action / streamed output};
\node[box,fill=green!6] (t) at (9.2,-.35) {Tool service\\ticket, order, refund};
\node[box,fill=orange!8] (d) at (14.1,1.15) {Specialist Agent\\optional delegation};
\draw[<->,thick] (u)--node[above,yshift=15pt]{client link C}(a);
\draw[<->,thick] (a.north east)--node[above,sloped]{model link M}(m.west);
\draw[<->,thick] (a.south east)--node[below,sloped]{tool link T}(t.west);
\draw[<->,thick] (a.east)--node[above]{coordination A}(d.west);
\node[note] at (2.0,-.35) {\textbf{Transit observer}\\lengths, gaps, connections;\\destinations when available};
\node[note] at (13.7,-.35) {\textbf{Recipient / active service}\\received arguments;\\optionally selected responses};
\node[draw=black!45,fill=black!3,align=left,text width=15.7cm,inner sep=6pt] at (7.9,-1.65) {
\textbf{Two protection questions:} Where does tool feedback appear next? What work does privacy cover execute before it is discarded? The original task and each observable service remain part of the protection decision.};
\end{tikzpicture}
\caption{Running support-ticket example and observation rights. A tool result returns to the Agent and can enter the next model request. Each encrypted link exposes its own traffic shape; a recipient additionally sees the request addressed to it. Concurrent tickets may share a connection on either link. The arrows represent application communication, not a one-to-one mapping from tool events to packets.}
\label{fig:chain}
\end{figure*}


\textbf{Defense placement determines both remaining exposure and permitted execution effects.} A tool-side intervention covers a particular interface. Its evaluation must also account for other interfaces available to the attacker and for any work the intervention causes. This SoK connects fifteen direct studies, forty-four settings and twenty-two established traffic-analysis anchors around these two requirements.

\paragraph{Follow information across services}
WRA recovers research queries from tool destinations, while ANT and the cloud benchmark recognize task labels on model connections~\cite{jeong2026wra,li2026ant,pouryousef2026traffic}. A contacted tool has yet another view: its addressed arguments. These observations change the unit of protection from a single transformed trace to a declared set of accessible interfaces. The field map identifies those interfaces, the rights needed to observe them and the evidence for each signal. Our forwarding pilot separately tests the incremental effect of supplementary feedback; its corrected results leave that causal question open.

\paragraph{Follow cover into execution}
Transport cover can terminate without invoking application work; an Agent decoy may execute queries or update state. WRA's published research-decoy comparison reduces query recovery, while its original-task utility scores concern a different visibility intervention~\cite{jeong2026wra}. The practical consequence is to measure privacy, original-task outcomes and resource use under the same cover policy. We illustrate this record with six paired policies and two fixed-weight planners over MCP Streamable HTTP. Deliberately adverse writes test whether a target-task predicate captures collateral changes; read, discarded and isolated cover provide the accompanying controls. This is an evaluation example, not a finding of harmful writes in WRA.

Website fingerprinting (WF) and machine-learning security already establish realistic populations, segmentation, transfer and adaptive evaluation~\cite{juarez2014critical,wang2016realistic,dyer2012peekaboo,mathews2023critical,arp2022dos,tramer2020adaptive}. Our contribution is their execution-aware systematization: a deployment and transfer map; two source-grounded syntheses connecting defense placement to exposure and task effects; and a worked procedure for reporting both together. The maps distinguish transit, recipient and active access instead of combining their scores. Source-dependency analysis identifies which comparisons survive without preprints. Together, these artifacts turn a defense's placement into concrete observation and outcome requirements.
